%% ma: L11-B15-0091
%% lop: 11
%% chuong: 5
%% bai: 15
%% dang: 11.15.5
%% loai: TLN
%% muc: VD
%% dap_an: "64"
%% nguon: "tai-lieu-hoc-tap-mon-toan-lop-11-theo-chuong-trinh-giao-duc-pho-thong-2018.pdf (Lớp Toán Thầy Duy, Phạm Lê Duy, THPT Chu Văn An), Chương 3 Bài 1, Phần 3, Câu 13 (tr.343)"
%% kien_thuc: "Lớp 11 Bài 15 (giới hạn của dãy số)"
%% trang_thai: cho-duyet
%% ly_do: "Dạng mới đề xuất (11.15.x), chờ giáo viên duyệt dạng; tài liệu gốc không có lời giải mẫu, đáp án do kho tự giải và kiểm bằng sympy."
\begin{ex}
Biết rằng giới hạn \[\lim\left(\sqrt{n^2-4n+5}+3-n\right)=\lim\left(\dfrac{5-a\cdot n}{\sqrt{n^2-4n+5}+n}+b\right)\] với $a;b$ là các số tự nhiên. Tính giá trị $S=a^b$.
\shortans{64}
\loigiai{\sloppy $\sqrt{n^2-4n+5}-n=\dfrac{5-4n}{\sqrt{n^2-4n+5}+n}$, nên $\sqrt{n^2-4n+5}+3-n=\dfrac{5-4n}{\sqrt{n^2-4n+5}+n}+3$. Vậy $a=4$, $b=3$ và $S=4^3=64$.}
\end{ex}
